% https://www.bilibili.com/video/BV1UugkzkEUE/?spm_id_from=333.1387.upload.video_card.click

\begin{lstlisting}[style=pythonstyle,numbers=left]
#!/usr/bin/env python3
"""
高级Python屏幕录制工具 - 支持区域选择、音频录制、格式选择、暂停恢复等
功能特色:
- 🎯 精确区域选择和全屏录制
- ⏸️ 暂停和恢复录制功能
- 🎬 多种视频格式和编码器支持
- ⚙️ 可调节帧率和质量
- 📁 自定义输出目录和文件名
- ⌨️ 快捷键支持 (Ctrl+Shift+S停止, Ctrl+Shift+P暂停)
- 📊 实时录制统计显示

作者: GitHub Copilot
版本: 2.1 (添加暂停/恢复功能)
"""

import cv2
import numpy as np
import tkinter as tk
from tkinter import ttk, messagebox, filedialog
import threading
import time
import os
import sys
from PIL import Image, ImageTk
import subprocess
import platform
import json
from datetime import datetime

class AdvancedScreenRecorder:
    def __init__(self):
        self.root = tk.Tk()
        self.root.title("屏幕录制工具 v2.1 (支持暂停/恢复)")
        self.root.geometry("700x600")
        self.root.resizable(True, True)
        
        # 录制参数
        self.is_recording = False
        self.is_paused = False  # 新增：暂停状态
        self.output_path = "recordings"
        self.fps = 30
        self.quality = 85
        self.format = "mp4"
        self.codec = "mp4v"
        
        # 录制区域 (x, y, width, height)
        self.record_area = None
        
        # 录制统计
        self.start_time = None
        self.pause_time = None  # 新增：暂停开始时间
        self.total_pause_duration = 0  # 新增：总暂停时长
        self.frame_count = 0
        self.video_writer = None
        self.record_thread = None  # 新增：录制线程引用
        
        # 创建输出目录
        os.makedirs(self.output_path, exist_ok=True)
        
        # 加载配置
        self.load_config()
        
        self.setup_ui()
        
        # 设置键盘快捷键
        self.setup_hotkeys()
        
        # 保存配置在退出时
        self.root.protocol("WM_DELETE_WINDOW", self.on_closing)
        
    def setup_ui(self):
        """设置用户界面"""
        # 创建主框架和滚动条
        main_canvas = tk.Canvas(self.root)
        scrollbar = ttk.Scrollbar(self.root, orient="vertical", command=main_canvas.yview)
        scrollable_frame = ttk.Frame(main_canvas)
        
        scrollable_frame.bind(
            "<Configure>",
            lambda e: main_canvas.configure(scrollregion=main_canvas.bbox("all"))
        )
        
        main_canvas.create_window((0, 0), window=scrollable_frame, anchor="nw")
        main_canvas.configure(yscrollcommand=scrollbar.set)
        
        main_frame = ttk.Frame(scrollable_frame, padding="15")
        main_frame.grid(row=0, column=0, sticky=(tk.W, tk.E, tk.N, tk.S))
        
        # 标题
        title_label = ttk.Label(main_frame, text="屏幕录制工具 v2.1", 
                               font=("Arial", 18, "bold"))
        title_label.grid(row=0, column=0, columnspan=3, pady=(0, 20))
        
        # 录制区域设置
        area_frame = ttk.LabelFrame(main_frame, text="录制区域设置", padding="10")
        area_frame.grid(row=1, column=0, columnspan=3, sticky=(tk.W, tk.E), pady=(0, 15))
        
        ttk.Button(area_frame, text="🎯 选择录制区域", 
                  command=self.select_area).grid(row=0, column=0, padx=(0, 10), sticky=tk.W)
        
        ttk.Button(area_frame, text="🖥️ 全屏录制", 
                  command=self.select_fullscreen).grid(row=0, column=1, padx=(0, 10), sticky=tk.W)
        
        ttk.Button(area_frame, text="🪟 选择窗口", 
                  command=self.select_window).grid(row=0, column=2, sticky=tk.W)
        
        self.area_label = ttk.Label(area_frame, text="❌ 未选择区域", foreground="red")
        self.area_label.grid(row=1, column=0, columnspan=3, sticky=tk.W, pady=(10, 0))
        
        # 录制设置
        settings_frame = ttk.LabelFrame(main_frame, text="录制设置", padding="10")
        settings_frame.grid(row=2, column=0, columnspan=3, sticky=(tk.W, tk.E), pady=(0, 15))
        
        # 第一行：FPS和质量
        ttk.Label(settings_frame, text="帧率(FPS):").grid(row=0, column=0, sticky=tk.W)
        self.fps_var = tk.StringVar(value=str(self.fps))
        fps_spinbox = ttk.Spinbox(settings_frame, from_=10, to=60, 
                                 textvariable=self.fps_var, width=8)
        fps_spinbox.grid(row=0, column=1, padx=(5, 15), sticky=tk.W)
        
        ttk.Label(settings_frame, text="质量(1-100):").grid(row=0, column=2, sticky=tk.W)
        self.quality_var = tk.StringVar(value=str(self.quality))
        quality_spinbox = ttk.Spinbox(settings_frame, from_=10, to=100, 
                                     textvariable=self.quality_var, width=8)
        quality_spinbox.grid(row=0, column=3, padx=(5, 0), sticky=tk.W)
        
        # 第二行：格式和编码器
        ttk.Label(settings_frame, text="输出格式:").grid(row=1, column=0, sticky=tk.W, pady=(10, 0))
        self.format_var = tk.StringVar(value=self.format)
        format_combo = ttk.Combobox(settings_frame, textvariable=self.format_var, 
                                   values=["mp4", "avi", "mov", "mkv"], width=8, state="readonly")
        format_combo.grid(row=1, column=1, padx=(5, 15), sticky=tk.W, pady=(10, 0))
        
        ttk.Label(settings_frame, text="编码器:").grid(row=1, column=2, sticky=tk.W, pady=(10, 0))
        self.codec_var = tk.StringVar(value=self.codec)
        codec_combo = ttk.Combobox(settings_frame, textvariable=self.codec_var,
                                  values=["mp4v", "XVID", "MJPG", "X264"], width=8, state="readonly")
        codec_combo.grid(row=1, column=3, padx=(5, 0), sticky=tk.W, pady=(10, 0))
        
        # 输出设置
        output_frame = ttk.LabelFrame(main_frame, text="输出设置", padding="10")
        output_frame.grid(row=3, column=0, columnspan=3, sticky=(tk.W, tk.E), pady=(0, 15))
        
        ttk.Label(output_frame, text="保存目录:").grid(row=0, column=0, sticky=tk.W)
        ttk.Button(output_frame, text="📁 选择目录", 
                  command=self.select_output_dir).grid(row=0, column=1, padx=(10, 0))
        
        self.output_label = ttk.Label(output_frame, text=f"📂 当前: {self.output_path}")
        self.output_label.grid(row=1, column=0, columnspan=2, sticky=tk.W, pady=(5, 0))
        
        # 文件名设置
        ttk.Label(output_frame, text="文件名前缀:").grid(row=2, column=0, sticky=tk.W, pady=(10, 0))
        self.filename_var = tk.StringVar(value="screen_record")
        filename_entry = ttk.Entry(output_frame, textvariable=self.filename_var, width=20)
        filename_entry.grid(row=2, column=1, padx=(10, 0), sticky=tk.W, pady=(10, 0))
        
        # 高级选项
        advanced_frame = ttk.LabelFrame(main_frame, text="高级选项", padding="10")
        advanced_frame.grid(row=4, column=0, columnspan=3, sticky=(tk.W, tk.E), pady=(0, 15))
        
        # 录制选项
        self.show_cursor_var = tk.BooleanVar(value=True)
        ttk.Checkbutton(advanced_frame, text="显示鼠标光标", 
                       variable=self.show_cursor_var).grid(row=0, column=0, sticky=tk.W)
        
        self.minimize_on_record_var = tk.BooleanVar(value=True)
        ttk.Checkbutton(advanced_frame, text="录制时最小化窗口", 
                       variable=self.minimize_on_record_var).grid(row=0, column=1, sticky=tk.W, padx=(20, 0))
        
        # 热键设置
        ttk.Label(advanced_frame, text="热键设置:").grid(row=1, column=0, sticky=tk.W, pady=(10, 0))
        hotkey_info = ttk.Label(advanced_frame, text="停止: Ctrl+Shift+S | 暂停: Ctrl+Shift+P", foreground="blue")
        hotkey_info.grid(row=1, column=1, sticky=tk.W, padx=(20, 0), pady=(10, 0))
        
        # 控制按钮
        control_frame = ttk.Frame(main_frame)
        control_frame.grid(row=5, column=0, columnspan=3, pady=(20, 0))
        
        self.record_button = ttk.Button(control_frame, text="🔴 开始录制", 
                                       command=self.toggle_recording)
        self.record_button.grid(row=0, column=0, padx=(0, 15))
        
        self.pause_button = ttk.Button(control_frame, text="⏸️ 暂停", 
                                      command=self.toggle_pause, state="disabled")
        self.pause_button.grid(row=0, column=1, padx=(0, 15))
        
        ttk.Button(control_frame, text="📁 打开录制目录", 
                  command=self.open_output_dir).grid(row=0, column=2, padx=(0, 15))
        
        ttk.Button(control_frame, text="⚙️ 测试录制", 
                  command=self.test_recording).grid(row=0, column=3, padx=(0, 15))
        
        ttk.Button(control_frame, text="❌ 退出", 
                  command=self.on_closing).grid(row=0, column=4)
        
        # 状态和统计显示
        status_frame = ttk.LabelFrame(main_frame, text="录制状态", padding="10")
        status_frame.grid(row=6, column=0, columnspan=3, sticky=(tk.W, tk.E), pady=(20, 0))
        
        self.status_label = ttk.Label(status_frame, text="🟢 就绪", 
                                     font=("Arial", 12, "bold"))
        self.status_label.grid(row=0, column=0, sticky=tk.W)
        
        self.stats_label = ttk.Label(status_frame, text="", 
                                    font=("Arial", 10))
        self.stats_label.grid(row=1, column=0, sticky=tk.W, pady=(5, 0))
        
        # 进度条
        self.progress = ttk.Progressbar(status_frame, mode='indeterminate')
        self.progress.grid(row=2, column=0, sticky=(tk.W, tk.E), pady=(10, 0))
        
        # 配置网格权重
        for i in range(4):
            settings_frame.columnconfigure(i, weight=1)
        for i in range(2):
            output_frame.columnconfigure(i, weight=1)
            advanced_frame.columnconfigure(i, weight=1)
        main_frame.columnconfigure(2, weight=1)
        status_frame.columnconfigure(0, weight=1)
        
        # 配置滚动
        main_canvas.pack(side="left", fill="both", expand=True)
        scrollbar.pack(side="right", fill="y")

    def select_area(self):
        """选择录制区域"""
        if self.minimize_on_record_var.get():
            self.root.withdraw()  # 隐藏主窗口
        else:
            self.root.iconify()  # 最小化主窗口
        time.sleep(0.2)  # 等待窗口隐藏
        
        # 创建区域选择窗口
        selector = AreaSelector()
        area = selector.select_area()
        
        self.root.deiconify()  # 显示主窗口
        
        if area:
            self.record_area = area
            x, y, w, h = area
            self.area_label.config(text=f"✅ 区域: ({x},{y}) 尺寸: {w}×{h}", foreground="green")
            self.status_label.config(text="✅ 已选择录制区域")
        else:
            self.status_label.config(text="❌ 未选择区域")

    def select_fullscreen(self):
        """选择全屏录制"""
        try:
            import pyautogui
            screen_width, screen_height = pyautogui.size()
            self.record_area = (0, 0, screen_width, screen_height)
            self.area_label.config(text=f"✅ 全屏: {screen_width}×{screen_height}", foreground="green")
            self.status_label.config(text="✅ 已选择全屏录制")
        except ImportError:
            messagebox.showerror("错误", "需要安装pyautogui库:\npip install pyautogui")

    def select_window(self):
        """选择特定窗口录制"""
        messagebox.showinfo("提示", "窗口选择功能开发中...\n请使用区域选择或全屏录制")

    def select_output_dir(self):
        """选择输出目录"""
        directory = filedialog.askdirectory(initialdir=self.output_path)
        if directory:
            self.output_path = directory
            self.output_label.config(text=f"📂 当前: {self.output_path}")

    def setup_hotkeys(self):
        """设置全局快捷键"""
        try:
            # 绑定窗口快捷键（仅在窗口有焦点时有效）
            self.root.bind('<Control-Shift-S>', lambda e: self.toggle_recording())
            self.root.bind('<Control-Shift-P>', lambda e: self.toggle_pause())
            
            # 注意：全局快捷键需要额外的库如pynput，这里先实现窗口快捷键
            print("快捷键已设置:")
            print("  Ctrl+Shift+S: 开始/停止录制")
            print("  Ctrl+Shift+P: 暂停/恢复录制")
            print("  (需要窗口有焦点时才有效)")
            
        except Exception as e:
            print(f"设置快捷键失败: {e}")

    def open_output_dir(self):
        """打开输出目录"""
        if platform.system() == "Windows":
            os.startfile(self.output_path)
        elif platform.system() == "Darwin":  # macOS
            subprocess.run(["open", self.output_path])
        else:  # Linux
            subprocess.run(["xdg-open", self.output_path])

    def test_recording(self):
        """测试录制功能"""
        if not self.record_area:
            messagebox.showerror("错误", "请先选择录制区域!")
            return
        
        try:
            import pyautogui
            # 测试截图
            x, y, width, height = self.record_area
            screenshot = pyautogui.screenshot(region=(x, y, width, height))
            
            # 保存测试截图
            test_file = os.path.join(self.output_path, "test_screenshot.png")
            screenshot.save(test_file)
            
            messagebox.showinfo("测试成功", f"录制区域测试通过！\n测试截图已保存: {test_file}")
            
        except Exception as e:
            messagebox.showerror("测试失败", f"录制测试失败: {str(e)}")

    def toggle_recording(self):
        """开始/停止录制"""
        if not self.is_recording:
            self.start_recording()
        else:
            self.stop_recording()

    def toggle_pause(self):
        """暂停/恢复录制"""
        if not self.is_recording:
            return
            
        if self.is_paused:
            self.resume_recording()
        else:
            self.pause_recording()

    def pause_recording(self):
        """暂停录制"""
        if not self.is_recording or self.is_paused:
            return
            
        self.is_paused = True
        self.pause_time = time.time()
        
        # 更新UI
        self.pause_button.config(text="▶️ 恢复")
        self.status_label.config(text="⏸️ 录制已暂停")
        self.progress.stop()
        
        # 显示窗口（如果需要）
        if self.minimize_on_record_var.get():
            self.root.deiconify()
            self.root.lift()

    def resume_recording(self):
        """恢复录制"""
        if not self.is_recording or not self.is_paused:
            return
            
        # 计算暂停时长
        if self.pause_time:
            self.total_pause_duration += time.time() - self.pause_time
            self.pause_time = None
            
        self.is_paused = False
        
        # 更新UI
        self.pause_button.config(text="⏸️ 暂停")
        self.status_label.config(text="🔴 正在录制...")
        self.progress.start(10)
        
        # 最小化窗口（如果需要）
        if self.minimize_on_record_var.get():
            self.root.iconify()

    def start_recording(self):
        """开始录制"""
        if not self.record_area:
            messagebox.showerror("错误", "请先选择录制区域!")
            return
        
        try:
            self.fps = int(self.fps_var.get())
            self.quality = int(self.quality_var.get())
            self.format = self.format_var.get()
            self.codec = self.codec_var.get()
        except ValueError:
            messagebox.showerror("错误", "请输入有效的数值!")
            return
        
        # 重置状态
        self.is_recording = True
        self.is_paused = False
        self.total_pause_duration = 0
        self.pause_time = None
        
        # 更新UI
        self.record_button.config(text="⏹️ 停止录制")
        self.pause_button.config(text="⏸️ 暂停", state="normal")
        self.status_label.config(text="🔴 正在录制...")
        self.progress.start(10)
        
        if self.minimize_on_record_var.get():
            self.root.iconify()  # 最小化窗口
        
        # 在新线程中开始录制
        self.record_thread = threading.Thread(target=self.record_screen)
        self.record_thread.daemon = True
        self.record_thread.start()

    def stop_recording(self):
        """停止录制"""
        self.is_recording = False
        self.is_paused = False
        
        # 更新UI
        self.record_button.config(text="🔴 开始录制")
        self.pause_button.config(text="⏸️ 暂停", state="disabled")
        self.progress.stop()
        
        # 显示窗口
        self.root.deiconify()
        self.root.lift()
        
        # 等待录制线程结束
        if self.record_thread and self.record_thread.is_alive():
            self.record_thread.join(timeout=2.0)

    def record_screen(self):
        """录制屏幕的主要逻辑"""
        try:
            import pyautogui
            
            # 禁用pyautogui的安全特性
            pyautogui.FAILSAFE = False
            
            x, y, width, height = self.record_area
            
            # 生成输出文件名
            timestamp = datetime.now().strftime("%Y%m%d_%H%M%S")
            filename = f"{self.filename_var.get()}_{timestamp}.{self.format}"
            output_file = os.path.join(self.output_path, filename)
            
            # 创建视频写入器
            fourcc = cv2.VideoWriter_fourcc(*self.codec)
            self.video_writer = cv2.VideoWriter(output_file, fourcc, self.fps, (width, height))
            
            if not self.video_writer.isOpened():
                raise Exception(f"无法创建视频文件: {output_file}")
            
            self.frame_count = 0
            self.start_time = time.time()
            frame_interval = 1.0 / self.fps
            
            while self.is_recording:
                # 如果暂停，等待恢复
                while self.is_paused and self.is_recording:
                    time.sleep(0.1)
                    continue
                
                # 如果录制已停止，退出循环
                if not self.is_recording:
                    break
                    
                frame_start = time.time()
                
                # 截取屏幕
                screenshot = pyautogui.screenshot(region=(x, y, width, height))
                frame = np.array(screenshot)
                
                # 转换颜色格式 (PIL使用RGB，OpenCV使用BGR)
                frame = cv2.cvtColor(frame, cv2.COLOR_RGB2BGR)
                
                # 写入帧
                self.video_writer.write(frame)
                self.frame_count += 1
                
                # 更新状态（计算实际录制时间，排除暂停时间）
                current_time = time.time()
                elapsed = current_time - self.start_time - self.total_pause_duration
                fps_actual = self.frame_count / elapsed if elapsed > 0 else 0
                
                # 根据暂停状态显示不同信息
                if self.is_paused:
                    status_text = f"帧数: {self.frame_count} | 录制时间: {elapsed:.1f}s | 实际FPS: {fps_actual:.1f} | ⏸️ 已暂停"
                else:
                    status_text = f"帧数: {self.frame_count} | 录制时间: {elapsed:.1f}s | 实际FPS: {fps_actual:.1f}"
                
                self.root.after(0, lambda text=status_text: self.stats_label.config(text=text))
                
                # 控制帧率
                frame_time = time.time() - frame_start
                sleep_time = max(0, frame_interval - frame_time)
                if sleep_time > 0:
                    time.sleep(sleep_time)
            
            # 关闭视频写入器
            if self.video_writer:
                self.video_writer.release()
                self.video_writer = None
            
            # 更新UI
            final_time = time.time() - self.start_time - self.total_pause_duration
            final_fps = self.frame_count / final_time if final_time > 0 else 0
            
            self.root.after(0, lambda: self.status_label.config(
                text=f"✅ 录制完成! 文件: {filename}"))
            self.root.after(0, lambda: self.stats_label.config(
                text=f"总帧数: {self.frame_count} | 录制时间: {final_time:.1f}s | 平均FPS: {final_fps:.1f} | 暂停时长: {self.total_pause_duration:.1f}s"))
            
            # 保存配置
            self.save_config()
            
        except Exception as e:
            self.root.after(0, lambda: messagebox.showerror("录制错误", str(e)))
            self.root.after(0, lambda: self.status_label.config(text="❌ 录制失败"))
        finally:
            self.is_recording = False
            self.is_paused = False
            if self.video_writer:
                self.video_writer.release()
                self.video_writer = None
            self.root.after(0, lambda: self.record_button.config(text="🔴 开始录制"))
            self.root.after(0, lambda: self.pause_button.config(text="⏸️ 暂停", state="disabled"))
            self.root.after(0, lambda: self.progress.stop())

    def load_config(self):
        """加载配置"""
        config_file = os.path.join(self.output_path, "recorder_config.json")
        try:
            if os.path.exists(config_file):
                with open(config_file, 'r', encoding='utf-8') as f:
                    config = json.load(f)
                    self.fps = config.get('fps', 30)
                    self.quality = config.get('quality', 85)
                    self.format = config.get('format', 'mp4')
                    self.codec = config.get('codec', 'mp4v')
                    self.output_path = config.get('output_path', self.output_path)
        except Exception as e:
            print(f"加载配置失败: {e}")

    def save_config(self):
        """保存配置"""
        config_file = os.path.join(self.output_path, "recorder_config.json")
        try:
            config = {
                'fps': self.fps,
                'quality': self.quality,
                'format': self.format,
                'codec': self.codec,
                'output_path': self.output_path
            }
            with open(config_file, 'w', encoding='utf-8') as f:
                json.dump(config, f, indent=2, ensure_ascii=False)
        except Exception as e:
            print(f"保存配置失败: {e}")

    def on_closing(self):
        """关闭程序时的处理"""
        if self.is_recording:
            result = messagebox.askyesno("确认", "正在录制中，确定要退出吗？\n录制将被停止并保存当前进度。")
            if result:
                self.stop_recording()
                # 等待录制线程结束
                if self.record_thread and self.record_thread.is_alive():
                    self.record_thread.join(timeout=3.0)
                time.sleep(0.5)  # 等待录制停止
            else:
                return
        
        self.save_config()
        self.root.destroy()

class AreaSelector:
    """改进的区域选择器"""
    
    def __init__(self):
        self.start_x = None
        self.start_y = None
        self.end_x = None
        self.end_y = None
        self.selected_area = None
        self.rect_id = None
        
    def select_area(self):
        """选择屏幕区域"""
        try:
            import pyautogui
            
            # 创建全屏透明窗口
            self.overlay = tk.Toplevel()
            self.overlay.attributes('-fullscreen', True)
            self.overlay.attributes('-alpha', 0.3)
            self.overlay.configure(bg='black')
            self.overlay.attributes('-topmost', True)
            
            # 获取屏幕尺寸
            screen_width = self.overlay.winfo_screenwidth()
            screen_height = self.overlay.winfo_screenheight()
            
            # 创建画布
            self.canvas = tk.Canvas(self.overlay, width=screen_width, height=screen_height,
                                   bg='black', highlightthickness=0, cursor='crosshair')
            self.canvas.pack()
            
            # 添加说明文字
            help_text = "🎯 拖拽鼠标选择录制区域 | ESC取消 | Enter确认"
            self.canvas.create_text(screen_width//2, 50, 
                                   text=help_text,
                                   fill="white", font=("Arial", 16, "bold"))
            
            # 绑定事件
            self.canvas.bind("<Button-1>", self.on_click)
            self.canvas.bind("<B1-Motion>", self.on_drag)
            self.canvas.bind("<ButtonRelease-1>", self.on_release)
            self.overlay.bind("<Escape>", self.cancel_selection)
            self.overlay.bind("<Return>", self.confirm_selection)
            self.overlay.bind("<KeyPress>", self.on_key)
            
            self.overlay.focus_set()
            
            # 等待选择完成
            self.overlay.wait_window()
            
            return self.selected_area
            
        except ImportError:
            messagebox.showerror("错误", "需要安装pyautogui库:\npip install pyautogui")
            return None
    
    def on_click(self, event):
        """鼠标按下"""
        self.start_x = event.x
        self.start_y = event.y
        # 清除之前的选择
        if self.rect_id:
            self.canvas.delete(self.rect_id)
    
    def on_drag(self, event):
        """鼠标拖拽"""
        if self.start_x is not None:
            # 清除之前的矩形
            if self.rect_id:
                self.canvas.delete(self.rect_id)
            
            # 绘制新矩形
            self.rect_id = self.canvas.create_rectangle(
                self.start_x, self.start_y, event.x, event.y,
                outline="red", width=3, fill="red", stipple="gray25")
            
            # 显示尺寸信息
            width = abs(event.x - self.start_x)
            height = abs(event.y - self.start_y)
            self.canvas.create_text(
                (self.start_x + event.x) // 2, 
                (self.start_y + event.y) // 2,
                text=f"{width} × {height}",
                fill="yellow", font=("Arial", 14, "bold"), tags="size_text")
    
    def on_release(self, event):
        """鼠标释放"""
        if self.start_x is not None:
            self.end_x = event.x
            self.end_y = event.y
            self.update_selection()
    
    def on_key(self, event):
        """键盘事件"""
        if event.keysym == "Return":
            self.confirm_selection(event)
        elif event.keysym == "Escape":
            self.cancel_selection(event)
    
    def update_selection(self):
        """更新选择区域"""
        if self.start_x is not None and self.end_x is not None:
            # 计算选择区域
            x1, y1 = min(self.start_x, self.end_x), min(self.start_y, self.end_y)
            x2, y2 = max(self.start_x, self.end_x), max(self.start_y, self.end_y)
            
            width = x2 - x1
            height = y2 - y1
            
            if width > 20 and height > 20:  # 最小尺寸检查
                self.selected_area = (x1, y1, width, height)
                # 更新显示
                self.canvas.delete("size_text")
                self.canvas.create_text(
                    x1 + width // 2, y1 + height // 2,
                    text=f"✅ {width} × {height}\n按Enter确认",
                    fill="lime", font=("Arial", 14, "bold"), tags="size_text")
    
    def confirm_selection(self, event):
        """确认选择"""
        if self.selected_area:
            self.overlay.destroy()
        else:
            self.canvas.create_text(
                self.overlay.winfo_screenwidth() // 2, 100,
                text="❌ 请先选择一个区域！",
                fill="red", font=("Arial", 16, "bold"))
    
    def cancel_selection(self, event):
        """取消选择"""
        self.selected_area = None
        self.overlay.destroy()

def check_dependencies():
    """检查依赖库"""
    required_packages = [
        ("cv2", "opencv-python"),
        ("PIL", "Pillow"),
        ("numpy", "numpy"),
        ("pyautogui", "pyautogui")
    ]
    
    missing_packages = []
    
    for package, install_name in required_packages:
        try:
            __import__(package)
        except ImportError:
            missing_packages.append(install_name)
    
    if missing_packages:
        message = "缺少以下依赖库，请安装:\n\n"
        message += "pip install " + " ".join(missing_packages)
        messagebox.showerror("依赖检查", message)
        return False
    
    return True

def main():
    """主函数"""
    if not check_dependencies():
        return
    
    app = AdvancedScreenRecorder()
    
    try:
        app.root.mainloop()
    except KeyboardInterrupt:
        print("\n程序被用户中断")
    except Exception as e:
        messagebox.showerror("错误", f"程序运行出错: {str(e)}")

if __name__ == "__main__":
    main()
\end{lstlisting}
